\icode{E-EXM-FI-L2-260625-001}
逐差法的优点是提高实验数据的利用率，减小了随机误差的影响，另外也可减小了实验中仪器误差分量。因此逐差法是一种常用的数据处理方法，不要把逐差法局限在求加速度这一种情景中，如

	某同学用力的传感器研究弹簧弹力与形变量关系时，在弹性限度内，每次使弹簧伸长一定长度并记录相应的力传感器示数，如下表所示

	\begin{table}[htbp]
	\centering
	\begin{tblr}{|c|c|c|c|c|}
	\hline
	弹簧总长度/\,\qty{}{cm}&12.00&14.00&16.00&18.00\\
	\hline
	\(F\)/\,\qty{}{N}&30.00&31.04&32.02&33.02\\
	\hline
	\end{tblr}
	\end{table}
	根据表格，要求尽可能多的利用测量数据，计算弹簧的劲度系数\(k =\) \underline{\hbox to 4em{}}\,\qty{}{N/m}（结果保留三位有效数字）。